\section{Discusión}
Las figuras evidencian impactos esperados y focos de trabajo. La mejora en tiempos deriva de flujos digitalizados y validaciones; la adopción depende de la estrategia de cambio y una UX consistente; los defectos iniciales guían esfuerzos de QA; el SLA proyectado es factible con monitoreo/escalado; y el ROI apoya la inversión. Limitaciones: los datos son estimaciones que deberán confirmarse en piloto.